\section{Lecture 16: Training Deep Networks}

\subsection{Motivation}

In the previous lecture, we introduced backpropagation as an efficient method for computing gradients.

\medskip

\noindent
However, in practice, training deep neural networks presents significant challenges.

\medskip

\noindent
\textbf{Key question.}  
Why is training deep networks difficult, and how can we address these challenges?

\medskip

\noindent
\textbf{Guiding idea.}  
Training behavior depends on initialization, gradient dynamics, and optimization structure.

\subsection{Initialization Strategies}

Neural networks are typically initialized randomly:
\[
W_\ell \sim \text{random distribution}.
\]

\medskip

\noindent
\textbf{Why initialization matters.}

\begin{itemize}
    \item Affects scale of activations
    \item Influences gradient flow
    \item Determines training stability
\end{itemize}

\medskip

\noindent
\textbf{Poor initialization.}

\begin{itemize}
    \item Too small:
    \begin{itemize}
        \item Signals vanish across layers
    \end{itemize}

    \item Too large:
    \begin{itemize}
        \item Activations and gradients explode
    \end{itemize}
\end{itemize}

\medskip

\noindent
\textbf{Principle of variance preservation.}

\begin{itemize}
    \item Choose initialization so that:
    \begin{itemize}
        \item activations maintain stable variance across layers
        \item gradients do not shrink or explode
    \end{itemize}
\end{itemize}

\medskip

\noindent
\textbf{Insight.}  
Initialization sets the starting geometry of optimization.

\subsection{Vanishing and Exploding Gradients}

During backpropagation, gradients are multiplied repeatedly:

\[
\frac{\partial L}{\partial x} =
\prod_{\ell} W_\ell^\top \cdot \sigma'(z_\ell).
\]

\medskip

\noindent
\textbf{Vanishing gradients.}
\begin{itemize}
    \item Gradients shrink exponentially
    \item Early layers learn very slowly
\end{itemize}

\medskip

\noindent
\textbf{Exploding gradients.}
\begin{itemize}
    \item Gradients grow exponentially
    \item Training becomes unstable
\end{itemize}

\medskip

\noindent
\textbf{Causes.}
\begin{itemize}
    \item Repeated matrix multiplications
    \item Activation functions with small or large derivatives
\end{itemize}

\medskip

\noindent
\textbf{Consequences.}
\begin{itemize}
    \item Difficulty training deep networks
    \item Sensitivity to hyperparameters
\end{itemize}

\medskip

\noindent
\textbf{Insight.}  
Gradient flow degrades as depth increases.

\subsection{Optimization Challenges}

Deep learning optimization differs from classical settings.

\medskip

\noindent
\textbf{Non-convexity.}
\begin{itemize}
    \item Many local minima and saddle points
\end{itemize}

\medskip

\noindent
\textbf{Ill-conditioning.}
\begin{itemize}
    \item Different directions have different curvature
    \item Slows convergence
\end{itemize}

\medskip

\noindent
\textbf{Coupled parameters.}
\begin{itemize}
    \item Layers interact in complex ways
\end{itemize}

\medskip

\noindent
\textbf{Empirical observation.}
\begin{itemize}
    \item Gradient-based methods often work well despite theoretical difficulty
\end{itemize}

\medskip

\noindent
\textbf{Insight.}  
Optimization landscapes are complex but structured.

\subsection{Empirical Training Behavior}

In practice, several patterns emerge:

\medskip

\noindent
\textbf{Rapid initial progress.}
\begin{itemize}
    \item Loss decreases quickly early in training
\end{itemize}

\medskip

\noindent
\textbf{Plateaus.}
\begin{itemize}
    \item Slow progress due to flat regions or saddle points
\end{itemize}

\medskip

\noindent
\textbf{Generalization behavior.}
\begin{itemize}
    \item Training loss may continue decreasing
    \item Test performance may peak earlier (overfitting)
\end{itemize}

\medskip

\noindent
\textbf{Role of SGD.}
\begin{itemize}
    \item Noise helps escape poor regions
    \item Encourages flatter solutions
\end{itemize}

\medskip

\noindent
\textbf{Insight.}  
Training dynamics are influenced by both optimization and stochasticity.

\subsection{Practical Remedies (Preview)}

Several techniques address training challenges:

\begin{itemize}
    \item Careful initialization
    \item Choice of activation functions
    \item Normalization methods
    \item Architectural design (e.g., residual connections)
\end{itemize}

\medskip

\noindent
These will be studied in subsequent lectures.

\subsection{A Unifying Perspective}

Training deep networks involves:

\begin{itemize}
    \item \textbf{Initialization:} setting the starting point
    \item \textbf{Gradient flow:} propagation of learning signals
    \item \textbf{Optimization dynamics:} navigating complex landscapes
    \item \textbf{Stochasticity:} shaping generalization
\end{itemize}

\medskip

\noindent
\textbf{Core idea.}  
Effective training requires controlling both signal propagation and optimization dynamics.

\medskip

\noindent
\textbf{Key insight.}  
Depth introduces both expressive power and optimization difficulty.

\subsection{Outlook}

In this lecture, we examined:

\begin{itemize}
    \item initialization strategies,
    \item vanishing and exploding gradients,
    \item optimization challenges,
    \item empirical training behavior.
\end{itemize}

\medskip

\noindent
To address these issues, we introduce techniques that stabilize training.

\medskip

\noindent
In the next lecture, we study activation functions and architectural design choices that improve performance.

\medskip

\noindent
\textbf{Guiding principle.}  
Successful deep learning requires both expressive models and stable training dynamics.